
\فصل{پیش‌زمینه و پژوهش‌های پیشین}

در این فصل به بررسی مبانی و کارهای پیشین مرتبط با موضوع پژوهش پرداخته می‌شود. از آنجا که درک درست روش پیشنهادی نیازمند آشنایی با برخی مفاهیم پایه‌ای در معماری‌های نوین شبکه‌های عصبی و سازوکارهای یادگیری است، در ابتدا به معرفی پیش‌نیازهای لازم مانند شبکه‌های {مبدل بینایی\پانویس{Vision Transformer}}\cite{vit} و انواع روش‌های {نرمال‌سازی\پانویس{Normalization}} لایه‌ها پرداخته خواهد شد. این مبانی نقش مهمی در طراحی مدل‌های مقاوم در برابر تغییر توزیع داده ایفا می‌کنند و در بسیاری از رویکردهای جدید انطباق زمان آزمون مورد استفاده قرار گرفته‌اند.

پس از آن، تمرکز اصلی این فصل بر بررسی خانواده‌ی روش‌های انطباق زمان آزمون است. این روش‌ها به طور ویژه برای مقابله با انباشت خطا و افت عملکرد مدل‌های یادگیری عمیق در شرایطی که توزیع داده‌های آزمون با داده‌های آموزش تفاوت دارد طراحی شده‌اند. در این بخش، به معرفی و دسته‌بندی رویکردهای موجود پرداخته می‌شود، به‌ویژه آن دسته از روش‌هایی که برای حالتهای پیوسته و جریان داده توسعه یافته‌اند.

در پایان، جمع‌بندی کلی از نقاط قوت و ضعف دسته‌های مختلف روش‌های انطباق زمان آزمون ارائه خواهد شد تا مسیر روشن‌تری برای معرفی معماری پیشنهادی در فصل بعد فراهم شود.

\قسمت{پیش‌زمینه‌}

برای درک بهتر روش پیشنهادی و تحلیل دقیق‌تر مسئله‌ی انطباق زمان آزمون، لازم است برخی مفاهیم پایه‌ای در حوزه‌ی شبکه‌های عصبی عمیق مرور شوند. در این بخش ابتدا شبکه‌های مبدل بینایی که یکی از معماری‌های مدرن و قدرتمند در بینایی ماشین محسوب می‌شوند معرفی می‌گردند. سپس، به بررسی مسئله‌ی نرمال‌سازی در شبکه‌های عصبی و انواع رویکردهای متداول آن پرداخته می‌شود. هدف از این مرور، ارائه‌ی دیدی روشن نسبت به اجزای کلیدی است که در بسیاری از روش‌های انطباق زمان آزمون مورد استفاده قرار می‌گیرند.

\زیرقسمت{مبدل بینایی}

معماری‌های مبتنی بر {مبدل\پانویس{Transformer}}\cite{vit} در ابتدا برای پردازش زبان طبیعی طراحی شدند و به دلیل توانایی در مدل‌سازی وابستگی‌های طولانی‌، موفقیت چشمگیری به دست آوردند. این موفقیت باعث شد تا پژوهشگران از ایده‌ی مبدل‌ها در بینایی ماشین نیز بهره بگیرند. شبکه‌ی مبدل بینایی یکی از نخستین معماری‌هایی است که به صورت مستقیم تصاویر را به قطعه‌های کوچک تقسیم کرده و هر قطعه را مشابه یک {تکواژه\پانویس{Token}} در مبدل متنی پردازش می‌کند.

مزیت اصلی مبدل بینایی در مقایسه با شبکه‌های {کانولوشنی\fontencoding{CNN}}، توانایی آن در یادگیری روابط غیرمحلی و انعطاف در مقیاس‌پذیری است. این ویژگی‌ها موجب شده مبدل بینایی به یکی از معماری‌های پایه‌ای در روش‌های مدرن انطباق زمان آزمون تبدیل شود. ایده‌ی اصلی آن است که یک تصویر ورودی 
\(
x \in \mathbb{R}^{H \times W \times C}
\)
را به قطعات کوچکی به نام {قطعه\پانویس{Patch}} تقسیم کرده و هر قطعه را مشابه یک تکواژه متنی در نظر بگیرد.  

ابتدا تصویر به \(N\) قطعه با ابعاد \(P \times P\) تقسیم می‌شود:  
\[
N = \frac{HW}{P^2}.
\]

هر قطعه به صورت یک بردار ستونی درآورده شده و توسط یک لایه‌ی خطی \lr{(Patch Embedding)} به فضای ویژگی با بعد \(D\) نگاشت می‌شود:  
\[
z_0 = [x_p^1 \mathbf{E}; \; x_p^2 \mathbf{E}; \dots; x_p^N \mathbf{E}] + \mathbf{E}_{pos},
\]  
که در آن 
\(\mathbf{E} \in \mathbb{R}^{(P^2 C) \times D}\)
ماتریس وزنی لایه‌ی تعبیه و 
\(\mathbf{E}_{pos}\)
بردارهای تعبیه‌ی موقعیتی هستند که اطلاعات مکانی هر قطعه را حفظ می‌کنند.  

سپس توالی بردارهای تعبیه‌شده وارد لایه‌های مختلف مبدل شامل {سازوکار توجه چندسری\پانویس{Multi-Head AtTENTion}} و {شبکه‌های پیش‌خور\پانویس{Feed Forward}} می‌شوند. در نهایت، به هر دنباله‌ی تکواژه‌های حاصل از تعبیه‌ی قطعات، یک بردار ویژه موسوم به تکواژه‌ی \lr{[CLS]} افزوده می‌شود. این تکواژه در طول عبور از لایه‌های مبدل، اطلاعات کلی تمام قطعات تصویر را تجمیع کرده و به‌عنوان نماینده‌ی کل تصویر عمل می‌کند. خروجی نهایی متناظر با تکواژه‌ی \lr{[CLS]} پس از گذر از یک لایه‌ی پیش‌بینی معمولاً یک {سر چند‌لایه‌ای پرسپترون\پانویس{Multi-Layer Perceptron Head}} به‌عنوان نمایش کلی تصویر برای تخمین برچسب کلاس مورد استفاده قرار می‌گیرد.

ساختار کلی معماری مبدل بینایی در شکل \ref{fig:vit} نشان داده شده است.

\begin{figure}[]
	\centering
	\includegraphics[width=1\linewidth]{images/vit.png}
	\caption[معماری کلی مبدل بینایی]{معماری کلی مبدل بینایی\cite{vit}.}
	\label{fig:vit}
\end{figure}
\زیرقسمت{نرمال‌سازی لایه‌ها}

نرمال‌سازی یکی از مؤلفه‌های کلیدی شبکه‌های عصبی عمیق است که نقش مهمی در پایداری آموزش، جلوگیری از {محو شدن\پانویس{Vanishing Gradient}} یا {انفجار گرادیان‌ها\پانویس{Exploding Gradient}} و بهبود کارایی مدل دارد. در ادامه، چند روش پرکاربرد نرمال‌سازی معرفی می‌شوند.  

\زیرزیرقسمت{نرمال‌سازی دسته‌ای}  

در {نرمال‌سازی دسته‌ای\پانویس{Batch Normalization}}\cite{bn}، برای هر ویژگی در طول {میان‌دسته\پانویس{Mini-batch}} میانگین و واریانس محاسبه شده و نرمال‌سازی طبق رابطه‌ی \ref{eq:bn} صورت می‌گیرد\cite{bn}: 
\begin{equation}
	\hat{x}_i = \frac{x_i - \mu_\mathcal{B}}{\sqrt{\sigma_\mathcal{B}^2 + \epsilon}}, \quad
	y_i = \gamma \hat{x}_i + \beta
	\label{eq:bn}
\end{equation}

که در آن \(\mu_\mathcal{B}\) و \(\sigma_\mathcal{B}^2\) میانگین و واریانس روی {دسته\پانویس{batch}}، و \(\gamma, \beta\) پارامترهای قابل یادگیری هستند.  

\زیرزیرقسمت{نرمال‌سازی لایه‌ای}  

در {نرمال‌سازی لایه‌ای\پانویس{Layer Normalization}}\cite{ln} نرمال‌سازی روی تمام ویژگی‌های یک نمونه به صورت مستقل طبق رابطه‌ی \ref{eq:ln} انجام می‌شود\cite{ln}:
\begin{equation}
	\hat{x}_i = \frac{x_i - \mu_{L}}{\sqrt{\sigma_{L}^2 + \epsilon}}
	\label{eq:ln}
\end{equation}

که در آن \(\mu_{L}\) و \(\sigma_{L}^2\) به ترتیب میانگین و واریانس ویژگی‌های همان نمونه هستند. این روش به دلیل استقلال از دسته، انتخاب اصلی در معماری‌های مبتنی بر مبدل از جمله مبذل بینایی است.  

\زیرزیرقسمت{نرمال‌سازی نمونه‌ای}  

در {نرمال‌سازی نمونه‌ای\پانویس{Instance Normalization}}\cite{in} برای هر نمونه و هر کانال به صورت جداگانه نرمال‌سازی طبق رابطه‌ی \ref{eq:in} انجام می‌شود\cite{in}:
\begin{equation}
	\hat{x}_{i,c} = \frac{x_{i,c} - \mu_{i,c}}{\sqrt{\sigma_{i,c}^2 + \epsilon}}
	\label{eq:in}
\end{equation}

که در آن \(\mu_{i,c}\) و \(\sigma_{i,c}^2\) به ترتیب میانگین و واریانس روی کانال \(c\) در نمونه‌ی \(i\) هستند.  

\زیرزیرقسمت{نرمال‌سازی گروهی}  

در {نرمال‌سازی گروهی\پانویس{Group Normalization}}\cite{gn} ویژگی‌ها به چند گروه تقسیم شده و نرمال‌سازی در هر گروه به طور مستقل طبق رابطه‌ی \ref{eq:gn} انجام می‌شود\cite{gn}:
\begin{equation}
	\hat{x}_i = \frac{x_i - \mu_{G}}{\sqrt{\sigma_{G}^2 + \epsilon}}
	\label{eq:gn}
\end{equation}

که در آن \(\mu_G, \sigma_G^2\) میانگین و واریانس محاسبه‌شده در هر گروه از ویژگی‌ها هستند.  

به طور کلی، در مسئله‌ی انطباق زمان آزمون که داده‌ها به صورت جریانی و با اندازه‌ی دسته کوچک یا حتی واحد ارائه می‌شوند، روش‌های نرمال‌سازی لایه‌ای و گروهی به دلیل استقلال از اندازه‌ی دسته مناسب‌تر از نرمال سازی دسته‌ای هستند. این موضوع در طراحی بسیاری از الگوریتم‌های مدرن انطباق زمان آزمون مورد توجه قرار گرفته است\cite{sar}.

\قسمت{مروری بر روش‌های انطباق در زمان آزمون برخط}

مسئله‌ی {انطباق پیوسته در زمان آزمون\پانویس{Online Test-Time Adaptation (OTTA)}} به شرایطی اشاره دارد که در آن یک مدل از پیش آموزش‌دیده بر روی حوزه‌ی منبع \(D_S\) با استفاده از یک دسته‌بند \(f_S\) باید به‌صورت پیوسته و مستمر، برچسب نمونه‌های ورودی را در جریان داده‌های جدید \( \{B_1, B_2, \dots \} \) پیش‌بینی کند، به‌گونه‌ای که دانش کسب‌شده از دسته‌های قبلی بتواند برای انطباق با دسته‌ی فعلی مورد استفاده قرار گیرد. 

در این حالت، هر دسته‌ ممکن است از توزیع متفاوتی نسبت به داده‌های منبع نمونه‌برداری شده باشد، که چالش‌هایی نظیر انباشت خطا و فراموشی فاجعه‌بار را ایجاد می‌کند. به‌طور‌کلی روش‌های انطباق پیوسته در زمان آزمون را می‌توان بر اساس روش‌های اصلی به چند دسته به صورت زیر تقسیم کرد:

\begin{itemize}
	\item {\textbf{روش‌های مبتنی بر آنتروپی:}\پانویس{Entropy-based Methods}} شامل روش‌هایی مانند \lr{TENT}\cite{TENT} و \lr{SAR}\cite{sar} و \lr{DeYo}\cite{deyo} که هدفشان کاهش عدم قطعیت پیش‌بینی‌ها در زمان آزمون است.
	\item {\textbf{روش‌های مبتنی بر انرژی:}\پانویس{Energy-based Methods}} شامل \lr{TEA}\cite{tea} و \lr{AEA}\cite{aea} که از توابع انرژی برای هم‌ترازی و پایدارسازی انطباق استفاده می‌کنند.
	\item {\textbf{برچسب‌زنی خودکار:}\پانویس{Pseudo-labeling}} روش‌هایی مانند \lr{T3A}\cite{t3a} و \lr{SHOT}\cite{shot} که از پیش‌بینی‌های مدل برای تولید برچسب‌های موقت استفاده می‌کنند و مدل را بازآموزی می‌کنند.
	\item {\textbf{تنظیم مبتنی بر سازگاری:}\پانویس{Consistency Regularization}} شامل \lr{REM}\cite{rem} و \lr{PETAL}\cite{petal} که با استفاده از معیارهای سازگاری، یادگیری پایدار را تضمین می‌کنند.
	\item {\textbf{تنظیم ضد فراموشی:}\پانویس{Anti-forgetting Regularization}} مانند \lr{CoTTA}\cite{cotta} و \lr{EATA}\cite{eata} که سازوکارهایی برای جلوگیری از فراموشی دانش پیشین ارائه می‌دهند.
\end{itemize}

در ادامه‌ی فصل، هر یک از این دسته‌ها به‌طور مفصل بررسی شده و مزایا و محدودیت‌های آن‌ها در حالتهای جریان داده و تغییر توزیع تحلیل می‌گردد. این بررسی، زمینه‌ی لازم برای ارائه‌ی روش پیشنهادی در فصل بعد را فراهم می‌کند.

\زیرقسمت{روش‌های مبتنی بر آنتروپی}

در این دسته از روش‌ها، ایده‌ی اصلی آن است که مدل در زمان آزمون باید خروجی‌هایی با قطعیت بالاتر تولید کند. 
برای حالتهای انطباق پیوسته، معمولاً برچسب‌های واقعی داده‌های آزمون در دسترس نیستند، 
بنابراین معیارهایی مانند آنتروپی توزیع احتمال خروجی به‌عنوان شاخصی برای طمینان مدل به کار گرفته می‌شوند. 

مبنای این رویکردها آن است که اگر مدل به‌خوبی با توزیع جدید منطبق شود، توزیع پیش‌بینی آن متمرکزتر یعنی با آنتروپی پایین‌تر خواهد بود. 
بر این اساس، انطباق مدل از طریق به‌روزرسانی پارامترها در جهت کاهش آنتروپی خروجی صورت می‌گیرد. 
مزیت اصلی این دسته، سادگی در پیاده‌سازی و سربار محاسباتی ناچیز است. 
با این حال، روش‌های مبتنی بر آنتروپی نسبت به {اغتشاش‌\پانویس{noise}} و خطاهای تجمعی حساس هستند؛ زیرا کاهش آنتروپی در صورت وجود پیش‌بینی‌های اشتباه، 
می‌تواند باعث تقویت خطا و افت کارایی در بلندمدت شود. \\
این خانواده در عمل به‌عنوان یکی از پرکاربردترین رویکردهای انطباق زمان آزمون مطرح شده 
و الهام‌بخش بسیاری از توسعه‌های بعدی مانند روش‌های ضد فراموشی و مبتنی بر انرژی بوده است.

\زیرزیرقسمت{\lr{TENT}}

روش \lr{TENT}\پانویس{Fully Test-time Adaptation by Entropy Minimization} به‌عنوان یکی از اولین و پرکاربردترین رویکردهای مبتنی بر آنتروپی مطرح شد. 
ایده اصلی این روش از یک تحلیل تجربی آغاز می‌شود: نویسندگان ابتدا بر روی یک مجموعه‌داده مشخص، 
رابطه میان آنتروپی خروجی مدل و نرخ خطا را بررسی کردند. 
نتایج این تحلیل نشان داد که همبستگی بالایی میان این دو وجود دارد؛ 
به‌طور خاص، نمونه‌هایی که مدل برای آن‌ها خروجی با آنتروپی پایین تولید می‌کند، 
اغلب با دقت بیشتری طبقه‌بندی می‌شوند. 
به بیان دیگر، کاهش آنتروپی پیش‌بینی‌ها در یک بازه مشخص معادل کاهش نرخ خطا است. این مورد در شکل \ref{fig:TENT-fig1} آورده شده‌است.

\begin{figure}[]
	\centering
	\includegraphics[width=1\linewidth]{images/TENT-fig1.png}
	\caption[رابطه‌ی آنتروپی با نرخ خطا و تابع هزینه در روش \lr{TENT}]{رابطه‌ی آنتروپی با نرخ خطا و تابع هزینه در مجموعه داده \lr{CIFAR-100-C}. آنتروپی می‌تواند به‌عنوان شاخص اطمینان مدل در زمان آزمون استفاده شود \cite{TENT}.}
	\label{fig:TENT-fig1}
\end{figure}

با تکیه بر این مشاهده، \lr{TENT} پیشنهاد می‌دهد که در زمان آزمون می‌توان 
تابع هدفی مبتنی بر آنتروپی تعریف کرده و آن را در جریان داده‌های بدون برچسب کمینه کرد. 
هدف مسئله طبق رابطه‌ی \ref{eq:TENT} به‌صورت زیر تعریف می‌شود\cite{tent}:  

\begin{equation}
	\min_{\theta} \; \mathbb{E}_{x \sim \mathcal{D}_{\text{test}}} 
	H(p_{\theta}(y|x)) 
	= \min_{\theta} \; \mathbb{E}_{x \sim \mathcal{D}_{\text{test}}} 
	\left[ - \sum_{c} p_{\theta}(y=c|x) \log p_{\theta}(y=c|x) \right]
	\label{eq:TENT}
\end{equation}

که در آن \(p_{\theta}(y|x)\) توزیع احتمال خروجی مدل برای نمونه آزمون \(x\) است 
و \(H(\cdot)\) آنتروپی را نشان می‌دهد. 

برای پیاده‌سازی عملی این ایده، \lr{TENT} تمامی پارامترهای شبکه را ثابت نگه می‌دارد 
و تنها وزن‌ها و میانگین/واریانس‌های لایه‌های نرمال‌سازی دسته‌ای را 
در طول فرآیند آزمون به‌روزرسانی می‌کند. 
این طراحی چند مزیت مهم دارد: 
\begin{itemize}
	\item \textbf{کارایی محاسباتی بالا:} به‌روزرسانی تعداد محدودی پارامتر موجب می‌شود که روش بسیار سبک بوده و قابل استفاده در حالتهای برخط و بلادرنگ باشد.  
	\item \textbf{نیاز نداشتن به داده برچسب‌دار:} کل فرآیند انطباق بر اساس داده‌های بدون برچسب آزمون انجام می‌شود.  
	\item \textbf{سادگی پیاده‌سازی:} این روش نیازی به تغییر معماری شبکه یا طراحی پیچیده ندارد.  
\end{itemize}

با وجود این، محدودیت مهم \lr{TENT} آن است که تمرکز صرف بر کاهش آنتروپی می‌تواند در شرایطی که داده‌ها دارای اغتشاش یا شامل نمونه‌های دشوار هستند، 
به تقویت پیش‌بینی‌های نادرست و انباشت خطا منجر شود. 
در نتیجه، در حالت‌های پیوسته ممکن است کارایی مدل کاهش یابد و مواردی چون انباشت خطا و فراموشی فاجعه‌بار به وجود بیایید. در واقع این روش بسیار مستعد {فروپاشی مدل\پانویس{Model Collapse}} و یافتن {راه‌حل‌های پیش‌پاافتاده\پانویس{Trivial Solutions}} می‌باشد.
این ضعف لزوم توسعه‌ی روش‌های هوشمندتر را نشان می‌دهد، به‌گونه‌ای که برای مقابله با این مشکلات از روش‌های ابتکارمند برای شناسایی و حذف نمونه‌های مستعد خطا می‌بایست استفاده شود.

\زیرزیرقسمت{\lr{SAR}}

روش \lr{SAR}\پانویس{Sharpness-aware and Reliable Entropy Minimization}به‌طور ویژه برای مقابله با حالتهایی طراحی شده است که در آن استفاده از روش‌های مبتنی بر کاهش صرف آنتروپی بدون اعمال کنترل‌های مناسب، منجر به راه‌حل‌های پیش‌پاافتاده و فروپاشی مدل می‌شود. در این رویکرد، نویسندگان ابتدا به تحلیل دلایل شکست و فروپاشی مدل در برخی شرایط خاص می‌پردازند. 

\begin{figure}[]
	\centering
	\includegraphics[width=1\linewidth]{images/sar-motivation.png}
	\caption[تحلیل نمونه‌های شکست در روش کاهش آنتروپی برخط]{تحلیل نمونه‌های شکست در روش کاهش آنتروپی برخط \lr{TENT}\cite{TENT}. 
		\lr{(a)}, \lr{(b)} پیش‌بینی‌های مدل را در طول انطباق برخط نشان می‌دهند. 
		\lr{(c)} چگونگی تغییر اندازه‌ی گرادیان‌ها با و بدون فروپاشی مدل را نمایش می‌دهد. 
		\lr{(d)} رابطه‌ی بین آنتروپی نمونه و اندازه‌ی گرادیان‌ها را بررسی می‌کند. 
		تمام آزمایش‌ها روی مجموعه داده \lr{ImageNet-C} با اغتشاش گاوسی و مدل \lr{ResNet50 (GN)} انجام شده است،
		و سطح بالاتر \lr{severity} نشان‌دهنده‌ی تغییر توزیع شدیدتر است \cite{sar}.}
	\label{fig:sar-motivation}
\end{figure}

تحلیل‌ها نشان می‌دهد که در گام‌های خاصی از فرایند به‌روز‌رسانی وزن‌های مدل برای انطباق در زمان آزمون اندازه‌ی گرادیان‌های نمونه‌ها به طور ناگهانی افزایش یافته و بلافاصله کاهش می‌یابد. بررسی رابطه‌ی بین اندازه‌ی گرادیان و آنتروپی نشان می‌دهد که بخش قابل توجهی از نمونه‌هایی که مستعد خطا هستند دارای آنتروپی زیاد هستند و گرادیان‌های بزرگی تولید می‌کنند. این نمونه‌ها در نواحی ۱ و ۲ بخش \lr{d} شکل \ref{fig:sar-motivation} قرار دارند. با اعمال {فیلترهای\پانویس{Filter}} اولیه، بخش زیادی از نمونه‌های مستعد خطا حذف می‌شوند، اما همچنان برخی نمونه‌ها در ناحیه ۴ بخش \lr{d} شکل \ref{fig:sar-motivation} وجود دارند که با وجود آنتروپی کم، گرادیان‌های بزرگی تولید می‌کنند.  

برای حذف یا تعدیل اثر این نمونه‌ها، \lr{SAR} پیشنهاد می‌کند از روش \lr{SAM}\cite{sam} استفاده شود، به گونه‌ای که بهینه ساز مدل نقاطی از فضای تابع هزینه‌ی آنتروپی را پیدا کند که در آن‌ها مقدار تابع هزینه کمترین تغییر را دارد یعنی {کمینه‌ی مسطح\پانویس{Flat Minimum}}. به بیان دیگر، \lr{SAM} سعی می‌کند مدل را به سمت نقاطی هدایت کند که در اطراف آن‌ها تغییرات کوچک در پارامترها، تغییرات ناچیزی در آنتروپی خروجی ایجاد کند، و به این ترتیب از تولید گرادیان‌های بزرگ ناشی از نمونه‌های پرت جلوگیری شود. این رویکرد موجب می‌شود انطباق زمان آزمون پایدارتر و مقاوم‌تر در برابر فروپاشی مدل باشد.

هدف بهینه‌سازی در روش \lr{SAR} به صورت زیر طبق رابطه‌ی \ref{eq:sar} تعریف می‌شود\cite{sar}:

\begin{equation}
	\min_{\theta} \; S(x_t) \cdot \max_{\|\epsilon\|_2 \leq \rho} 
	H\big(f_{\theta + \epsilon}(x_t)\big)
	\label{eq:sar}
\end{equation}

که در آن \(x_t\) نمونه‌ی آزمون در زمان \(t\) است، 
\(f_{\theta + \epsilon}(x_t)\) پیش‌بینی مدل با پارامترهای \(\theta\) به همراه اختلال \(\epsilon\) است، 
\(H(\cdot)\) آنتروپی توزیع پیش‌بینی را نشان می‌دهد، 
\(\epsilon\) بردار اختلال با محدودیت اندازه‌ی \(\|\epsilon\|_2 \leq \rho\) است و 
\(S(x_t) = \mathbf{1}\big(H(f_\theta(x_t)) < E_0\big)\) شاخصی برای انتخاب نمونه‌هایی با آنتروپی کمتر از آستانه \(E_0\) می‌باشد.

روش \lr{SAR} با این ترکیب تنها نمونه‌های پایدار و دارای آنتروپی کم را وارد فرایند انطباق می‌کند و بدین ترتیب در تلاش است تا از فروپاشی مدل جلوگیری کند. این کار باعث می‌شود کاهش آنتروپی به‌صورت کنترل‌شده انجام شود و مدل در طول انطباق برخط پایدار باقی بماند، بدون اینکه به راه‌حل‌های پیش‌پاافتاده که هدف مسئله نیستند برسد.\\
با این حال، یکی از نقاط ضعف روش \lr{SAR} افزایش تعداد فراپارامترها است. برای مثال، تعیین حد تغییرات مجاز برای اعمال آشفتگی در وزن‌های مدل، انتخاب آستانه‌ی آنتروپی برای فیلتر کردن نمونه‌های نامطمئن، و مشخص کردن لایه‌هایی که اجازه‌ی اعمال آشفتگی در آن‌ها وجود دارد، همگی با توجه به ماهیت بدون‌ناظر و برخط مسئله چالش‌برانگیز هستند. به طور خاص، لایه‌های انتهایی شبکه‌های عصبی معمولاً حساسیت بیشتری نسبت به تغییرات دارند \cite{mummadi2021confidence,choi2022improving}. به همین دلیل، خود روش \lr{SAR} بخشی از لایه‌های پایانی شبکه را {منجمد\پانویس{Freeze}} کرده و تنها لایه‌های ابتدایی را برای انطباق آزاد می‌گذارد. انتخاب عمق مناسب برای اعمال تغییرات می‌تواند تأثیر قابل‌توجهی بر عملکرد نهایی مدل داشته باشد.

\زیرزیرقسمت{\lr{DeYo}}

روش \lr{DeYO}\پانویس{Destroy Your Object} نیز مانند \lr{SAR} به دنبال طراحی روشی برای شناسایی نمونه‌های پرخطر و حذف آن‌ها از فرایند انطباق است.
ایده‌ی اصلی این روش آن است که اتکای صرف بر آنتروپی
به‌عنوان معیار اطمینان مدل در زمان آزمون می‌تواند در حضور ویژگی‌های غیرمرتبط یا عوامل مخدوش‌کننده\پانویس{Spuriosity} 
گمراه‌کننده باشد. 
برای رفع این مشکل، \lr{DeYO} علاوه بر آنتروپی، معیاری به نام 
\lr{PLPD}\پانویس{Pseudo-Label Probability Difference} 
را وارد فرایند انطباق می‌کند. 
در این روش ابتدا یک {تبدیل\پانویس{Transformation}} روی نمونه‌ی ورودی اعمال می‌شود و این کار به دنبال این است تا ساختار تصویر را دگرگون کند و برای نمونه‌های با اطمینان بالا انتظار می‌رود که \lr{PLPD} مقدار زیادی داشته باشد. اگر این اختلاف برای یک نمونه ای از یک آستانه‌ای بیشتر نباشد از فرایند انطباق حذف می‌شود. نمونه‌ای از روش‌های بررسی شده برای دگرگون‌سازی تصویر در شکل ‌\ref{fig:deyo} قابل مشاهده است.\\
مجموعه‌ی نمونه‌های انتخاب ‌شده توسط مدل به‌صورت زیر تعریف می‌شود:

\begin{figure}[]
	\centering
	\includegraphics[width=1\linewidth]{images/deyo_images.pdf}
	\caption[تبدیل‌های اعمالی در روش \lr{DeYo}]{نمونه‌ای از تبدیل‌های اعمالی بر روی نمونه‌ها در زمان آزمون در روش \lr{DeYo}\cite{deyo}.}
	\label{fig:deyo}
\end{figure}

\begin{equation}
	S_{\theta}(x) = \left\{\, x \;\middle|\; \mathrm{Ent}_{\theta}(x) < \tau_{\text{Ent}}, \; 
	\mathrm{PLPD}_{\theta}(x, x') > \tau_{\text{PLPD}} \right\},
	\label{eq:deyo-set}
\end{equation}

که در آن \(\mathrm{Ent}_{\theta}(x)\) آنتروپی پیش‌بینی مدل برای نمونه \(x\)، 
\(\tau_{\text{Ent}}\) آستانه‌ی آنتروپی، \(x'\) نسخه‌ی مخدوش‌شده\پانویس{Object-Destructive Transformation} 
از \(x\) و \(\mathrm{PLPD}_{\theta}(x, x')\) اختلاف احتمال برچسب کاذب\پانویس{Pseudo-Label} بین پیش‌بینی‌های مدل روی \(x\) و \(x'\) است.  

برای کنترل اهمیت نمونه‌ها، وزن‌دهی بر اساس رابطه‌ی زیر تعریف می‌شود\cite{deyo}:

\begin{equation}
	\alpha_{\theta}(x) = \frac{1}{\exp\{(\mathrm{Ent}_{\theta}(x) - E_0)\} + \exp\{-\mathrm{PLPD}_{\theta}(x, x')\}},
	\label{eq:deyo-alpha}
\end{equation}

و در نهایت، تابع هزینه‌ی روش به صورت زیر ارائه می‌شود:

\begin{equation}
	L_{\mathrm{DeYO}}(x; \theta) = \alpha_{\theta}(x) \cdot \mathbf{I}\{x \in S_{\theta}(x)\} \; \mathrm{Ent}_{\theta}(x),
	\label{eq:deyo-loss}
\end{equation}

که در آن \(\mathbf{I}\{\cdot\}\) عملگر شاخص است.  

به بیان دیگر، طبق رابطه‌ی \ref{eq:deyo-set} تنها نمونه‌هایی که هم دارای آنتروپی پایین و هم اختلاف معنی‌دار در \lr{PLPD} باشند وارد فرآیند انطباق می‌شوند. 
وزن‌دهی \(\alpha_{\theta}(x)\) در رابطه‌ی \ref{eq:deyo-alpha} اهمیت نسبی هر نمونه را مشخص می‌کند و در نهایت، 
تابع هزینه‌ی \ref{eq:deyo-loss} کاهش آنتروپی را برای نمونه‌های معتبر و پایدار هدایت می‌کند. 
این طراحی باعث می‌شود که انطباق مدل در زمان آزمون از فروپاشی مدل جلوگیری کرده و تمرکز آن بر روی نمونه‌های معنادار باقی بماند.\\
با این حال، همان‌طور که از روابط مشخص است، پیچیدگی محاسباتی روش افزایش یافته و انتخاب دقیق آستانه‌های \(\tau_{\text{Ent}}\) و \(\tau_{\text{PLPD}}\) 
از چالش‌های اصلی آن محسوب می‌شود.
در مواردی که یک مدل بر اساس ویژگی‌های محلی خاص پیش‌بینی‌های دقیقی ارائه دهد، این پیش‌بینی‌ها حتی پس از اعمال آشفتگی ناشی از تخریب تصویر بدون تغییر باقی می‌مانند و در نتیجه مقدار \lr{PLPD} پایین خواهد بود. بنابراین، شناسایی نمونه‌های مفید برای انطباق زمان آزمون صرفاً بر اساس راهبرد \lr{PLPD} بالا ممکن است چندان مؤثر نباشد. 

\زیرقسمت{روش‌های مبتنی بر انرژی}

یکی از محدودیت‌های استفاده از تابع آنتروپی در انطباق زمان آزمون آن است که مدل را ناخواسته وادار می‌کند نسبت به یک کلاس خاص متعهد شود. این مسئله در شرایطی که نمونه‌ی ورودی ذاتاً همراه با ابهام یا ناشناختگی است، می‌تواند منجر به پیش‌بینی‌های بیش‌ازحد مطمئن و در نهایت فروپاشی مدل گردد.  

برای رفع این مشکل، روش \lr{TEA}\cite{tea} پیشنهاد می‌کند به جای آنتروپی از انرژی به‌عنوان معیار جایگزین استفاده شود. در این دسته از روش‌ها نشان داده می‌شود که همان رابطه‌ی مشاهده‌شده میان نرخ خطا و آنتروپی، میان نرخ خطا و انرژی نیز برقرار است. بنابراین، به‌جای واداشتن مدل به انتخاب قطعی یک کلاس، انرژی به‌عنوان شاخصی نرم‌تر عمل کرده و اجازه می‌دهد نمونه‌هایی که ذاتاً نامطمئن هستند، همچنان مقادیر انرژی پایینی داشته باشند.\\
بدین ترتیب، استفاده از انرژی می‌تواند خطر {بیش‌اعتمادی مدل\پانویس{Overconfidence}} را کاهش داده و فرآیند انطباق برخط را پایدارتر سازد. 

\زیرزیرقسمت{\lr{TEA}}

روش \lr{TEA}\پانویس{Test-time Energy Adaptation} به‌عنوان یکی از نخستین رویکردهای مبتنی بر انرژی در انطباق زمان آزمون معرفی شده است. 
ایده‌ی اصلی این روش آن است که به جای کمینه‌سازی مستقیم آنتروپی پیش‌بینی‌ها، 
از انرژی به‌عنوان معیاری جایگزین برای سنجش میزان اطمینان مدل استفاده شود.
همانند TENT این روش نیز با مطالعه‌ی تجربی رابطه‌ی میان نرخ خطا با انرژی به این ایده‌‌‌ی کلی می‌رسد.
در واقع، کاهش انرژی نمونه‌ها در یک بازه‌ی مشخص معادل کاهش نرخ خطا است. این مورد در شکل \ref{fig:tea-fig1} آورده شده‌است.

\begin{figure}[]
	\centering
	\includegraphics[width=1\linewidth]{images/tea-fig1.png}
	\caption[رابطه‌ی انرژی با نرخ خطا در روش \lr{TEA}]{رابطه‌ی انرژی با نرخ خطا در روش ‌\lr{TEA} \cite{tea}.}
	\label{fig:tea-fig1}
\end{figure}

ابتدا انرژی هر نمونه بر اساس رابطه‌ی \ref{eq:tea} تعریف می‌شود\cite{tea}:  

\begin{equation}
	E_{\theta}(x) = - \log \sum_{y} \exp \big(f_{\theta}(x)[y]\big),
	\label{eq:tea}
\end{equation}

که در آن \(f_{\theta}(x)[y]\) خروجی خام (logit) مدل برای کلاس \(y\) است. 
بر این اساس، احتمال پیش‌بینی‌شده برای نمونه‌ی آزمون طبق رابطه‌ی \ref{eq:tea-prob} تعریف می‌شود\cite{tea}:  

\begin{equation}
	p_{\theta}(x_{\text{test}}) = 
	\frac{\exp \big(-E_{\theta}(x_{\text{test}})\big)}{Z_{\theta}}, 
	\quad
	Z_{\theta} = \sum_{y} \exp \big(f_{\theta}(x_{\text{test}})[y]\big),
	\label{eq:tea-prob}
\end{equation}

که در آن \(Z_{\theta}\) نقش تابع نرمال‌ساز راایفا می‌کند.\\
گرادیان لگاریتم احتمال نسبت به پارامترهای مدل طبق رابطه‌ی \ref{eq:tea-grad} به دست می‌آید\cite{tea}:  

\begin{equation}
	\frac{\partial \log p_{\theta}(x_{\text{test}})}{\partial \theta} 
	= \mathbb{E}_{\tilde{x} \sim p_{\theta}} \left[\frac{\partial E_{\theta}(\tilde{x})}{\partial \theta}\right] 
	- \frac{\partial E_{\theta}(x_{\text{test}})}{\partial \theta}.
	\label{eq:tea-grad}
\end{equation}

برای محاسبه‌ی عبارت اول سمت راست تساوی در رابطه‌ی \ref{eq:tea-grad}، 
روش \lr{TEA} از نمونه‌گیری به کمک الگوریتم \lr{SGLD}\پانویس{Stochastic Gradient Langevin Dynamics} \cite{sgld} که روشی برای تقریب زنجیره مارکوف مونت‌کارلو جهت نمونه‌گیری کارآمد است، بهره می‌گیرد.
فرآیند به‌روزرسانی در این نمونه‌گیری طبق رابطه‌ی \ref{eq:tea-sgld} انجام می‌شود\cite{tea}:  

\begin{equation}
	\tilde{x}_{t+1} = \tilde{x}_{t} - \frac{\alpha}{2} 
	\frac{\partial E_{\theta}(\tilde{x}_{t})}{\partial \tilde{x}_{t}} + \alpha \epsilon,
	\quad \epsilon \sim \mathcal{N}(0, I),
	\label{eq:tea-sgld}
\end{equation}

که در آن \(\alpha\) نرخ یادگیری نمونه‌گیری است. 
بدین ترتیب، \lr{SGLD} یک تقریب عملی از محاسبه‌ی \(Z_{\theta}\) و نمونه‌گیری از توزیع مدل فراهم می‌کند. 

در نهایت، تابع هدف \lr{TEA} با در نظر گرفتن تفاوت انرژی بین نمونه‌ی مشاهده‌شده و نمونه‌های بازنمونه‌گیری‌شده تعریف می‌شود:  

\begin{equation}
	\max_{\theta} \; \min_{\tilde{x}} \; 
	\big[E_{\theta}(\tilde{x}) - E_{\theta}(x_{\text{test}})\big].
	\label{eq:tea-loss}
\end{equation}

این فرمول‌بندی باعث می‌شود مدل انرژی نمونه‌های واقعی را کاهش داده و در عین حال نمونه‌های مصنوعی را با انرژی بالاتر متمایز کند.  

به طور خلاصه، روش \lr{TEA} نشان می‌دهد که انرژی نیز مانند آنتروپی با نرخ خطا همبستگی بالایی دارد. 
با این حال، برخلاف آنتروپی، انرژی لزوماً مدل را مجبور نمی‌کند که به یک کلاس خاص متعهد شود.
این ویژگی موجب کاهش خطر \پانویس{overconfidence}{بیش‌اعتمادی مدل} و جلوگیری از فروپاشی مدل می‌شود 
وانطباق برخط پایدارتری را فراهم می‌سازد.

\زیرزیرقسمت{\lr{AEA}}

روش \lr{AEA}\پانویس{Adaptive Energy Alignment for Accelerating Test-Time Adaptation} با هدف بهبود و تسریع در فرایند انطباق زمان آزمون پیشنهاد می‌کند انرژی و آنتروپی به طور همزمان در فرایند انطباق در نظر گرفته شوند.
ایده‌ی اصلی این روش آن است که کاهش صرف آنتروپی به تنهایی مانع کاهش انرژی شده و تضمینی برای هم‌ترازی انرژی بین دامنه‌ی منبع و هدف ندارد. 
این کار از ترکیب سه مؤلفه‌ی اصلی {EM\پانویس{Entropy Minimization}}، {SFEA\پانویس{Source-free Energy Alignment}}، {LCS\پانویس{Logit Cosine Similarity}} تشکیل شده است که هر یک به تفضیل در ادامه بررسی می‌شوند.

ابتدا تابع هزینه‌ی \lr{EM} برای نمونه‌ی $x$ از دامنه‌ی هدف $p_T$ به صورت زیر طبق رابطه‌ی \ref{eq:em-loss} تعریف می‌شود\cite{aea}:

\begin{equation}
	\mathcal{L}_{EM}(x; \theta) = H(x; \theta) 
	= -\sum_{j=1}^{K} p_{\theta}(y_j|x) \log p_{\theta}(y_j|x),
	\label{eq:em-loss}
\end{equation}
که می‌توان آن را بر اساس تعریف \lr{softmax} به شکل مبتنی بر انرژی طبق رابطه‌ی \ref{eq:em-energy} بازنویسی کرد\cite{aea}:

\begin{equation}
	H(x; \theta) = \sum_{j=1}^{K} p_{\theta}(y_j|x) E_{\theta}(x, y_j) - E_{\theta}(x),
	\label{eq:em-energy}
\end{equation}
که در آن $E_{\theta}(x,y) = -f_{\theta}(x)[y]$ انرژی نمونه برای کلاس $y$ است.

در معادله‌ی دوم، جمله‌ی اول $\sum_{j=1}^{K} p_{\theta}(y_j|x) E_{\theta}(x, y_j)$ مدل را به سمت کاهش انرژی نمونه‌ها متناسب با احتمال پیش‌بینی هر کلاس هدایت می‌کند. 
این امر باعث می‌شود پیش‌بینی‌های با اطمینان بالا، حتی مطمئن‌تر شوند و \lr{logit} هر کلاس متناسب با میزان اطمینان آن افزایش یابد. 
در مقابل، جمله‌ی دوم $E_{\theta}(x)$ که انرژی آزاد نامیده می‌شود، می‌تواند به عنوان یک جمله‌ی جریمه در نظر گرفته شود که مقیاس کلی خروجی شبکه برای تمام کلاس‌ها را کاهش می‌دهد.

بنابراین، از \lr{SFEA} برای محاسبه‌ی انرژی مرجع از نمونه‌های با انرژی پایین در یک دسته‌ی هدف طبق رابطه‌ی \ref{eq:sfea-ref-energy} استفاده می‌شود\cite{aea}:

\begin{equation}
	\hat{E}^s = \frac{1}{|\hat{B}|} \sum_{i \in \hat{B}} E_\theta(x_i), 
	\quad 
	\hat{B} = \{i : E(x_i) \le \delta_B\},
	\label{eq:sfea-ref-energy}
\end{equation}

و تابع هزینه‌ی آن برای هر نمونه‌ی هدف به صورت زیر طبق رابطه‌ی \ref{eq:sfea-loss} تعریف می‌شود\cite{aea}:

\begin{equation}
	\mathcal{L}_{SFEA}(x; \theta) = \log \big(1 + \exp(E_\theta(x) - \hat{E}^s)\big).
	\label{eq:sfea-loss}
\end{equation}

استفاده از انرژی به طور معمول، جهت حرکت انرژی برای هر کلاس را به‌طور مشخص تعیین نمی‌کند. 
برای این منظور، \lr{AEA} از اطلاعات ماتریس وزن دسته‌بند آموزش‌دیده بهره می‌گیرد تا جهت تراز انرژی هر نمونه‌ی هدف در فضای خروجی شبکه مشخص شود که در رابطه‌ی \ref{eq:lcs-sim} آورده شده‌است\cite{aea}:

\begin{equation}
	S(x; \theta) = \sigma_{\cos} \big[f_\theta(x), \bar{f}_{\hat{y}}\big], 
	\quad 
	\bar{f}_{\hat{y}} = W^T W_{ŷ},
	\label{eq:lcs-sim}
\end{equation}
و سپس تابع هزینه‌ی \lr{LCS} به صورت زیر طبق رابطه‌ی \ref{eq:lcs-loss} تعریف می‌شود\cite{aea}:

\begin{equation}
	\mathcal{L}_{LCS}(x; \theta) = w(x) \cdot (1 - S(x; \theta)) \cdot \mathbb{I}\{C(x; \theta) \ge C_0\},
	\label{eq:lcs-loss}
\end{equation}

در نهایت، تابع هزینه‌ی کلی روش مذکور هم به صورت زیر طبق رابطه‌ی \ref{eq:total-loss} تعریف می‌شود\cite{aea}:

\begin{equation}
	\mathcal{L}_{total} = \mathcal{L}_{EM} + \lambda_1 \cdot \mathcal{L}_{SFEA} + \lambda_2 \cdot \mathcal{L}_{LCS}.
	\label{eq:total-loss}
\end{equation}


روش \lr{AEA} با ترکیب انرژی و آنتروپی به دنبال تسریع فرآیند انطباق در زمان آزمون است و بین بیش‌اعتمادی و نبود قطعیت مورد دوم را ترجیح می‌دهد تا فرایند انطباق دارای استحکام و انطباق بیشتری باشد.
نتایج تجربی نشان می‌دهد که این روش در برخی مجموعه‌داده‌ها و مدل‌ها عملکرد بهتری ارائه می‌دهد. 
با این حال، مشکل اصلی \lr{AEA} مربوط به هم‌ترازی انرژی و آنتروپی است؛ اگر مدل از پیش آموزش‌دیده مقیاس متفاوتی برای انرژی و آنتروپی داشته باشد، پارامترهای $\lambda_1$ و $\lambda_2$ می‌توانند به تغییرات بسیار حساس شوند و تنظیم دقیق آن‌ها با توجه به ماهیت برخط و بدون ناظر بودن مسئله، یک چالش عملیاتی جدی محسوب می‌شود.

\زیرقسمت{روش‌های مبتنی بر برچسب‌زنی خودکار}

روش‌های مبتنی بر برچسب‌زنی خودکار یکی دیگر از خانواده‌های رایج در انطباق زمان آزمون هستند. 
ایده‌ی اصلی این روش‌ها آن است که مدل از پیش‌بینی‌های خود روی داده‌های بدون برچسب آزمون به عنوان برچسب موقت استفاده کند و سپس با استفاده از این برچسب‌ها به‌روزرسانی پارامترها انجام می‌دهد. 
به بیان دیگر، برخلاف روش‌های مبتنی بر آنتروپی که تنها تلاش می‌کنند پیش‌بینی‌ها با قطعیت بیشتری تولید شوند، 
روش‌های برچسب‌زنی خودکار از خود پیش‌بینی‌های مدل به‌عنوان منبع اطلاعات برای یادگیری استفاده می‌کنند. 
این ویژگی باعث می‌شود مدل بتواند به شکل مستقیم با نمونه‌های جدید تعامل کند و به تدریج دانش خود را نسبت به توزیع فعلی داده افزایش دهد.

یکی از مزایای کلیدی این روش‌ها، توانایی کنترل اثر نمونه‌های دارای اغتشاش است؛ 
ممکن است برخی نمونه‌ها با وجود آنتروپی پایین، در عمل موجب هدایت مدل به راه‌حل‌های پیش‌پاافتاده شوند، 
اما روش‌های برچسب‌زنی خودکار که به دنبال اعمال نوعی از حافظه برای تولید برچسب برای نمونه‌ها در شرایط برخط هستند، 
می‌توانند اثر این نمونه‌ها را در طول انطباق کنترل کنند و از فروپاشی مدل جلوگیری نمایند.

مزیت اصلی این دسته سادگی پیاده‌سازی و قابلیت سازگاری با مدل‌های از پیش آموزش‌دیده است و نیازی به تغییر معماری شبکه یا داده‌های برچسب‌دار جدید ندارد. 
با این حال، یکی از چالش‌های اصلی آن حساسیت به خطاهای تجمعی است؛ زیرا همچنان برچسب‌های نادرست می‌توانند باعث هدایت مدل به راه‌حل‌های نادرست شوند و کاهش عملکرد طولانی‌مدت را ایجاد کنند. 
در ادامه روش‌های \lr{T3A} و \lr{SHOT} که از جمله نمونه‌های شاخص این دسته هستند بررسی می‌شوند.

\زیرزیرقسمت{\lr{T3A}}

روش  \lr{T3A}\پانویس{Test-Time Classifier Adjustment Module for Model-Agnostic Domain Generalization}، برای بهبود انطباق مدل‌های آموزش‌دیده روی داده‌های منبع با توزیع هدف در زمان آزمون طراحی شده است.
به مانند کار‌های قبلی در این روش نیز شبکه‌ی عصبی به دو بخش تقسیم می‌شود: استخراج‌کننده‌ی ویژگی \(f_\theta\) و دسته‌بند خطی \(q_\omega\) برای لایه‌ی آخر، که \(\theta\) و \(\omega\) پارامترهای آن‌ها هستند.

برای یک نمونه‌ی آزمون \(x\)، پیش‌بینی مدل به صورت حداکثر احتمال تقریب‌زده به شکل زیر است:

\begin{equation}
	\hat{y} = \arg\max_{y_k} q_\omega(Y = y_k \mid f_\theta(x)) 
	= \frac{\exp(z \cdot \omega_k)}{\sum_j \exp(z \cdot \omega_j)},
	\label{eq:t3a-pred}
\end{equation}

که در آن \(z = f_\theta(x)\)، و \(\omega_k \in \mathbb{R}^{\text{zdim}}\) وزن‌های مربوط به کلاس \(k\) و \(\text{zdim}\) بعد ویژگی‌های استخراج‌شده توسط \(f_\theta\) است. 
هر \(\omega_k\) به‌عنوان الگوی نمایشی کلاس \(k\) عمل می‌کند و پیش‌بینی با اندازه‌گیری فاصله به وسیله‌ی ضرب داخلی بین الگو و نمایه‌ی نمونه‌ی ورودی انجام می‌شود. از آنجا که این الگوها در داده‌های منبع آموزش دیده‌اند، تضمینی برای عملکرد بهینه در توزیع هدف وجود ندارد.

\lr{T3A} این الگوها را در زمان آزمون تنظیم می‌کند. فرض کنید دسته‌ای از داده‌های آزمون \(x\) در زمان \(t\) از توزیع هدف \(P(X) \neq P_d(X)\) نمونه‌برداری شده‌اند. برای هر نمونه، ابتدا با استفاده از برچسب‌های موقت \(\hat{y}\) ورودی افزوده می‌شود. سپس {مجموعه‌ی پشتیبانی\پانویس{Set Support}} برای کلاس \(k\) به صورت زیر طبق رابطه‌ی \ref{eq:t3a-support-update} \cite{t3a}به‌روزرسانی می‌شود:

\begin{equation}
	S^k_t = 
	\begin{cases}
		S^k_{t-1} \cup \{f_\theta(x) / \|f_\theta(x)\|\}, & \text{if } \hat{y} = y_k,\\
		S^k_{t-1}, & \text{otherwise},
	\end{cases}
	\label{eq:t3a-support-update}
\end{equation}

که \(\|\cdot\|\) اندازه‌ی \(L_2\) بردار است و \(S^k_0 = \{\omega_k / \|\omega_k\|\}\). اگر دسته‌ای از نمونه‌ها موجود باشد، همین فرایند برای هر نمونه تکرار می‌شود. سپس پیش‌بینی نهایی با استفاده از میانگین الگوهای کلاس طبق رابطه‌ی \ref{eq:t3a-pred-adjusted} انجام می‌گیرد\cite{t3a}:

\begin{equation}
	\hat{y} = \arg\max_{y_k} \gamma_c(Y = y_k \mid f_\theta(x)) 
	= \frac{\exp(z \cdot c_k)}{\sum_j \exp(z \cdot c_j)},
	\label{eq:t3a-pred-adjusted}
\end{equation}

که \(c_k\) مرکز کلاس \(k\) است و از عناصر \(S^k_t\) به صورت زیر طبق رابطه‌ی \ref{eq:t3a-centroid} محاسبه می‌شود\cite{t3a}:

\begin{equation}
	c_k = \frac{1}{|S^k_t|} \sum_{z \in S^k_t} z.
	\label{eq:t3a-centroid}
\end{equation}

برای جلوگیری از استفاده از نمونه‌های نامطمئن با برچسب موقت اشتباه، \lr{T3A} از آنتروپی پیش‌بینی \(\hat{y}\) به عنوان معیار فیلتر طبق رابطه‌ی \ref{eq:t3a-entropy} استفاده می‌کند\cite{t3a}:

\begin{equation}
	H_\omega(\hat{Y} \mid z) = - \sum_k q_\omega(\hat{Y} = y_k \mid z) \log q_\omega(\hat{Y} = y_k \mid z),
	\label{eq:t3a-entropy}
\end{equation}

و فقط نمونه‌هایی که آنتروپی آن‌ها کمتر از یک آستانه \(\alpha_k\) است، در مجموعه پشتیبانی باقی می‌مانند.

بدین ترتیب روش  \lr{T3A} با استفاده از مجموعه پشتیبانی و فیلتر آنتروپی، الگوهای کلاس را به‌روزرسانی می‌کند تا پیش‌بینی‌ها در توزیع هدف پایدار باقی بمانند. این سازوکار اجازه می‌دهد مدل در طول انطباق برخط از تاثیر نمونه‌های با برچسب موقت نادرست جلوگیری کند و از فروپاشی مدل جلوگیری شود.
از جمله معایب \lr{T3A} نیز می‌توان به وابستگی به انتخاب درست آستانه‌ی آنتروپی \(\alpha_k\) و احتمال اختصاص اشتباه برچسب موقت اشاره کرد، که می‌تواند نویز به مجموعه پشتیبانی اضافه کرده و عملکرد مدل را مختل کند.

\زیرزیرقسمت{\lr{SHOT}}

روش \lr{SHOT}\پانویس{Source Hypothesis Transfer for Unsupervised Domain Adaptation} یکی از شناخته‌شده‌ترین رویکردها در خانواده‌ی برچسب‌زنی خودکار است که هدف آن انتقال {فرضیه‌ی منبع\پانویس{Source Hypothesis}} به حوزه‌ی هدف بدون نیاز به دسترسی مستقیم به داده‌های منبع می‌باشد.   
ایده‌ی اصلی این روش در انطباق زمان آزمون این است که بخش دسته‌بند مدل منبع ثابت نگه داشته شود و تنها بخش استخراج ویژگی $g_t$ در حوزه‌ی هدف به‌روزرسانی گردد. این به‌روزرسانی با استفاده‌ی ترکیبی از معیارهای مبتنی بر آنتروپی و تنوع خروجی انجام می‌شود.  

در ابتدا، دو تابع هزینه‌ی کلیدی تعریف می‌شوند:  
(1) آنتروپی شرطی که مدل را طبق رابطه‌ی \ref{eq:l_ent} به سمت پیش‌بینی‌های با قطعیت بیشتر هدایت می‌کند\cite{shot}،  
(2) تابع تنوع که طبق رابطه‌ی \ref{eq:l_div} مانع فروپاشی مدل و هم‌جهت شدن تمام پیش‌بینی‌ها با یک کلاس خاص می‌شود\cite{shot}.  

\begin{equation}
	\label{eq:l_ent}
	\mathcal{L}_{ent}(f_t; X_t) = - \mathbb{E}_{x_t \sim X_t} \sum_{k=1}^{K} \delta_k(f_t(x_t)) \log \delta_k(f_t(x_t)),
\end{equation}

\begin{equation}
	\label{eq:l_div}
	\mathcal{L}_{div}(f_t; X_t) = \sum_{k=1}^{K} \hat{p}_k \log \hat{p}_k = D_{KL}(\hat{p} \parallel \tfrac{1}{K}\mathbf{1}_K) - \log K,
\end{equation}

که در آن $\hat{p} = \mathbb{E}_{x_t \sim X_t}[\delta(f_t(x_t))]$ میانگین پیش‌بینی‌ها در کل داده‌های هدف است.  

با وجود این دو مؤلفه، مشکل اصلی همچنان باقی می‌ماند: برخی داده‌های هدف ممکن است به دلیل تغییر توزیع با فرضیه‌ی منبع به‌درستی هم‌راستا نشوند.  
برای کاهش این مشکل، \lr{SHOT} از یک روش {خودنظارتی\پانویس{Self Supervised}} برای تولید برچسب‌های موقتی بهره می‌گیرد.  

ابتدا، برای هر کلاس در حوزه‌ی هدف یک {مرکز\پانویس{centroid}} محاسبه می‌شود که توزیع داده‌ها را بهتر نمایندگی می‌کند. این کار طبق رابطه‌ی \ref{eq:init_centroid} انجام می‌گیرد\cite{shot}:  

\begin{equation}
	\label{eq:init_centroid}
	c^{(0)}_k = \frac{\sum_{x_t \in X_t} \delta_k(\hat{f}_t(x_t)) \, \hat{g}_t(x_t)}{\sum_{x_t \in X_t} \delta_k(\hat{f}_t(x_t))},
\end{equation}

که در آن $\hat{f}_t = \hat{g}_t \circ h_t$ فرضیه‌ی فعلی مدل است.  

سپس، برای هر نمونه‌ی بدون برچسب در حوزه‌ی هدف، برچسب موقتی آن با توجه به نزدیک‌ترین مرکز کلاس در فضای نمایش طبق رابطه‌ی \ref{eq:pseudo_label} تعیین می‌شود\cite{shot}:

\begin{equation}
	\label{eq:pseudo_label}
	\hat{y}_t = \arg \min_k D_f(\hat{g}_t(x_t), c^{(0)}_k),
\end{equation}

که $D_f(a,b)$ فاصله‌ی کسینوسی بین بردارهای $a$ و $b$ را نشان می‌دهد.  

با استفاده از این برچسب‌ها، مراکز کلاس دوباره طبق رابطه‌ی \ref{eq:update_centroid} محاسبه می‌شوند تا پایداری بیشتری حاصل شود\cite{shot}:  

\begin{equation}
	\label{eq:update_centroid}
	c^{(1)}_k = \frac{\sum_{x_t \in X_t} \mathbf{1}[\hat{y}_t = k] \, \hat{g}_t(x_t)}{\sum_{x_t \in X_t} \mathbf{1}[\hat{y}_t = k]}.
\end{equation}

این فرایند تنها یک بار به‌روزرسانی می‌شود که طبق نتایج تجربی، برچسب‌های موقتی نسبتاً پایدار و قابل‌اعتماد تولید می‌کند.  

در نهایت، کل تابع هدف روش \lr{SHOT} با ترکیب این سه مؤلفه (آنتروپی شرطی، تنوع و برچسب‌زنی خودکار) طبق رابطه‌ی \ref{eq:shot_final} تعریف می‌شود\cite{shot}:  

\begin{equation}
	\label{eq:shot_final}
	\mathcal{L}(g_t) = \mathcal{L}_{ent}(h_t \circ g_t; X_t) + \mathcal{L}_{div}(h_t \circ g_t; X_t) - 
	\beta \, \mathbb{E}_{(x_t,\hat{y}_t)\sim X_t \times \hat{Y}_t} \sum_{k=1}^{K} \mathbf{1}[k=\hat{y}_t] \log \delta_k(h_t(g_t(x_t))),
\end{equation}

که در آن $\beta > 0$ یک ابرپارامتر کنترلی برای تنظیم تعادل بین اجزای مختلف است.  

به این ترتیب، روش \lr{SHOT} توانسته است مزایای رویکردهای مبتنی بر آنتروپی را با یک سازوکار پایدارتر برچسب‌زنی خودکار ترکیب کند تا هم انطباق سریع و هم کاهش اثر برچسب‌های دارای اغتشاش را تضمین نماید. با این حال، باید توجه داشت که این روش اساساً برای مسائل برخط و جریان داده‌ها طراحی نشده است و راهکار مشخصی برای جلوگیری از فروپاشی مدل در شرایط خاص ارائه نمی‌کند. 
همچنین در حالاتی که نمونه‌ها دارای برچسب‌های {نامتوازن\پانویس{Imbalanced}} باشند، عبارت مربوط به تنوع در تابع هزینه یعنی رابطه‌ی \ref{eq:l_div} می‌تواند به عنوان مانعی برای انطباق درست عمل کند و این روش با فرض توازن برچسب‌های نمونه‌ها در یک دسته طراحی شده است.

\زیرقسمت{روش‌های مبتنی بر سازگاری}

دسته‌ی دیگری از رویکردهای مهم در انطباق زمان آزمون، 
روش‌های مبتنی بر سازگاری هستند. 
ایده‌ی اصلی در این روش‌ها آن است که پیش‌بینی مدل برای یک نمونه‌ی ورودی 
نباید با تغییرات جزئی و بی‌اهمیت در داده 
مانند اغتشاشات و انواع {تبدیلات داده‌ای\پانویس{Data Augmentation}}
به‌طور چشمگیری تغییر کند.
به عبارت دیگر، مدل باید پیش‌بینی‌های سازگار و پایدار نسبت به نماهای متفاوت یک نمونه که می‌توانند از طریق اغتشاشات تصادفی یا تبدیلات ساختاری بر داده ایجاد شوند ارائه کند.

این رویکرد برخلاف روش‌های مبتنی بر آنتروپی یا انرژی که 
مستقیماً یک معیار بهینه‌سازی را تعریف می‌کنند، 
و همچنین در مقابل روش‌های مبتنی بر برچسب‌زنی خودکار که 
از به کارگیری شبه‌برچسب‌ها برای هدایت آموزش بهره می‌برند، 
بر روی کاهش حساسیت مدل به اغتشاشات ورودی تمرکز دارد. 
در نتیجه، این روش‌ها می‌توانند تعمیم‌پذیری بهتری در شرایطی که داده‌های آزمون 
دارای تغییرات ساختاری هستند را تضمین کنند.  

از میان روش‌های مطرح در این دسته، می‌توان به \lr{PETAL} و \lr{REM} اشاره کرد 
که هر یک با طراحی متفاوت، تلاش می‌کنند معیارهای سازگاری را برای بهبود دقت و پایداری 
در زمان آزمون به کار گیرند. 

\زیرزیرقسمت{\lr{PETAL}}

روش \lr{PETAL}\پانویس{A Probabilistic Framework for Lifelong Test-Time Adaptation} بر پایه‌ی یک دیدگاه احتمالاتی نسبت به مسئله‌ی انطباق زمان آزمون بنا شده است. \lr{PETAL} فرض می‌کند که می‌توان مدل منبع را طبق رابطه‌ی \ref{eq:posterior-approx} به‌صورت یک توزیع پسین تقریبی \(q(\theta)\) بر روی وزن‌ها نمایش داد\cite{petal}:  

\begin{equation}
	q(\theta) \approx p(\theta | D), 
	\label{eq:posterior-approx}
\end{equation}

این نمایش باعث می‌شود که در زمان آزمون، دانش مدل منبع از طریق این توزیع حفظ گردد. سپس با جایگذاری این تقریب در تابع {لاگرانژ\پانویس{Lagrange}}، به رابطه‌ی \ref{eq:tta-posterior} برای چگالی پسین با استفاده از داده‌های آزمون بدون برچسب می‌رسیم\cite{petal}:  

\begin{equation}
	\log p(\theta | D, U) = \log q(\theta) - \bar{\lambda}\frac{1}{M}\sum_{m=1}^{M} H_{xe}(y_m', y_m | \bar{x}),
	\label{eq:tta-posterior}
\end{equation}

که در آن، عبارت آنتروپی متقابل \(H_{xe}\) نقش مهمی در کاهش خطاهای ناشی از برچسب‌زنی اشتباه ایفا می‌کند.  

هم‌چنین ‌برای جلوگیری از انباشت خطا و فراموشی فاجعه بار این روش از سازوکار بازگردانی تصادفی معرفی شده در \cite{cotta} استفاده می‌کند. در این روش، پس از هر به‌روزرسانی گرادیان وزن‌ها، بخشی از وزن‌ها به حالت اولیه‌ی منبع طبق رابطه‌ی \ref{eq:stochastic-restore} بازگردانده می‌شوند\cite{petal}: 

\begin{equation}
	\theta^{t+1}_f = m \odot \theta^0_f + (1 - m) \odot \theta^{t+1}_f, 
	\label{eq:stochastic-restore}
\end{equation}

که در آن \(m \sim Bernoulli(\rho)\) {پوششی\پانویس{Mask}} است که مشخص می‌کند کدام پارامترها بازگردانده شوند.  

برای بهبود این فرآیند، \lr{PETAL} از ماتریس اطلاعات فیشر (\lr{FIM}) \cite{ewc} طبق رابطه‌ی \ref{eq:fisher-restore} به‌عنوان معیاری برای سنجش اهمیت پارامترها استفاده می‌کند. در این حالت، تنها پارامترهای کم‌اهمیت به مقادیر اولیه‌ی منبع بازگردانده می‌شوند\cite{petal}.

\begin{equation}
	m_p = 
	\begin{cases}
		1 & \text{if } F_p < \gamma \\
		0 & \text{otherwise},
	\end{cases}
	\label{eq:fisher-restore}
\end{equation}

که در آن \(\gamma = quantile(F, \delta)\) آستانه‌ای تعیین‌شده بر اساس توزیع مقادیر فیشر است.  

به طور خلاصه این روش با استفاده از بازنمایی احتمالاتی مدل منبع و سازوکار بازگردانی پارامترها توانسته است اثر فراموشی فاجعه‌بار و انباشت خطا را کاهش دهد و انطباق برخط روی داده‌های بدون برچسب را پایدارتر کند. این روش نیاز به دسترسی به داده‌های منبع ندارد و دانش اولیه مدل را حفظ می‌کند. با این حال، پیچیدگی محاسباتی بالا، حساسیت به تنظیم فرا‌پارامتر‌ها و محدودیت در جریان داده‌های بسیار تغییرپذیر از جمله معایب آن محسوب می‌شوند. استفاده از تقریب {گوسی\پانویس{Gaussian}} برای توزیع پسین نیز ممکن است تمامی پیچیدگی فضای وزن‌ها را به‌طور کامل مدل‌سازی نکند.

\زیرزیرقسمت{روش \lr{REM}} 

روش \lr{REM}\پانویس{Ranked Entropy Minimization} بر پایه‌ی ایجاد زنجیره‌های پوشش\پانویس{Mask Chains} برای نمونه‌های آزمون توسعه یافته است. ایده‌ی اصلی این است که اعمال پوشش‌های متفاوت روی تصاویر نمونه باعث تغییر مستقیم آنتروپی پیش‌بینی می‌شود؛ به‌گونه‌ای که هرچه درصد پوشش بیشتر باشد، آنتروپی نمونه افزایش می‌یابد و در نتیجه دقت کاهش می‌یابد. برای مقابله با انباشت خطا و فروپاشی مدل، \lr{REM} پیشنهاد می‌دهد که به جای اعمال تابع هزینه‌ی آنتروپی، از تابع هزینه‌ی آنتروپی متقاطع\پانویس{Cross-Entropy} بین نمونه‌های متوالی در زنجیره استفاده شود. این کار باعث می‌شود تنها با چند {عبور جلو\پانویس{Forward Pass}} و بدون نیاز به به‌روزرسانی کل پارامترهای شبکه، انعطاف‌پذیری و استحکام مدل افزایش یابد.  

تابع هزینه‌ی \lr{MCL}\پانویس{Masked Consistency Loss} برای اطمینان از سازگاری پیش‌بینی‌ها در زنجیره‌های پوشش به صورت زیر طبق رابطه‌ی \ref{eq:rem-mcl} تعریف می‌شود\cite{rem}:  

\begin{equation}
	L_{\text{MCL}} = \sum_{i<j}^{MN} H(f_t(x_j), \text{sg}(f_t(x_i))),
	\label{eq:rem-mcl}
\end{equation}

که در آن \(H(p, q)\) آنتروپی متقابل بین دو توزیع احتمال \(p\) و \(q\) است، \(MN = \{0, m_1, m_2, \dots, m_N\}\) مجموعه‌ای از نسبت‌های پوشش است، و \(\text{sg}\) عملگر توقف گرادیان\پانویس{Stop-Gradient} را نشان می‌دهد.  

برای حفظ ترتیب آنتروپی و جلوگیری از پیش‌بینی‌های بیش‌اعتماد، \lr{REM} تابع \lr{ERL}\پانویس{Entropy Ranking Loss} را طبق رابطه‌ی \ref{eq:rem-erl} معرفی می‌کند\cite{rem}:  

\begin{equation}
	L_{\text{ERL}} = \sum_{i<j}^{MN} \max(0, S(f_t(x_i)) - \text{sg}(S(f_t(x_j))) + m),
	\label{eq:rem-erl}
\end{equation}

که \(m\) {حاشیه\پانویس{Margin}} است و هدف آن حفظ این اصل است که نمونه‌هایی با نسبت پوشش پایین‌تر، باید آنتروپی کمتری نسبت به نمونه‌های با نسبت پوشش بالاتر داشته باشند.
عبارت \lr{MCL} با کاهش اختلاف‌ها در توزیع‌های پیش‌بینی رتبه‌بندی‌شده، به طور غیرمستقیم آنتروپی پیش‌بینی‌ها را کاهش می‌دهد. با این حال، تفاوت‌های کوچک در توزیع پیش‌بینی می‌تواند باعث انطباق کند شود و همچنین ممکن است پیش‌بینی‌ها به سمت پس‌زمینه سوگیری داشته باشند، به ویژه زمانی که اجسام در تصویر مخفی یا نیمه‌پوشیده باشند. برای مقابله با این چالش و تکمیل سازوکار \lr{MCL}، تابع \lr{ERL} معرفی شده است که با حفظ ترتیب صریح آنتروپی نمونه‌ها با نسبت پوشش متفاوت، از پیش‌بینی‌های بیش‌اعتماد و سوگیری به سمت پس‌زمینه جلوگیری می‌کند و سازگاری سریع‌تر مدل را در طول زنجیره‌های پوشش فراهم می‌سازد.

تابع هزینه‌ی کل به صورت ترکیبی از \lr{MCL} و \lr{ERL}  طبق رابطه‌ی \ref{eq:rem_total} تعریف می‌شود\cite{rem}:  

\begin{equation}
	L_{\text{REM}} = L_{\text{MCL}} + \lambda \cdot L_{\text{ERL}},
	\label{eq:rem_total}
\end{equation}

عبارت \lr{MCL} با کاهش اختلاف‌ها در توزیع‌های پیش‌بینی رتبه‌بندی‌شده، به طور غیرمستقیم آنتروپی پیش‌بینی‌ها را کاهش می‌دهد. با این حال، تفاوت‌های کوچک در توزیع پیش‌بینی می‌تواند باعث انطباق کند شود و همچنین ممکن است پیش‌بینی‌ها به سمت پس‌زمینه سوگیری داشته باشند\پانویس{Biased}، به ویژه زمانی که اجسام در تصویر مخفی یا نیمه‌پوشیده باشند. برای مقابله با این چالش و تکمیل سازوکار \lr{MCL}، تابع \lr{ERL} معرفی شده است، که با حفظ ترتیب صریح آنتروپی نمونه‌ها با نسبت پوشش متفاوت، از پیش‌بینی‌های بیش‌اعتماد و سوگیری به سمت پس‌زمینه جلوگیری می‌کند و سازگاری سریع‌تر مدل را در طول زنجیره‌های پوشش فراهم می‌سازد.  

در کل، روش \lr{REM} یک رویکرد جدید برای مسئله‌ی انطباق برخط زمان آزمون ارائه می‌دهد و به طور مؤثر به حفظ استحکام و انعطاف ‌پذیری مدل کمک می‌کند. با این حال، شایان ذکر است که طول زنجیره و نسبت پوشش‌ها به عنوان فراپارامترهای کلیدی این روش مطرح هستند و برای دستیابی به بهترین عملکرد نیازمند تنظیم دقیق می‌باشند. مشابه سایر روش‌های مبتنی بر آنتروپی، \lr{REM} سازوکاری مشخص برای حذف نمونه‌های نامطمئن ارائه نمی‌دهد و این موضوع به ویژه در حالتهای انطباق برخط زمان آزمون می‌تواند چالش‌برانگیز باشد.  

\زیرقسمت{روش‌های تنظیم ضد‌ فراموشی}

همان‌طور که پیش‌تر اشاره شد، یکی از چالش‌های اساسی در انطباق برخط زمان آزمون، پدیده‌ی فراموشی فاجعه‌بار و انباشت خطا است. در این وضعیت، به‌روزرسانی‌های متوالی مدل بر اساس نمونه‌های جدید می‌تواند به از دست رفتن تدریجی دانش پیشین و در نتیجه کاهش عملکرد کلی منجر شود. این مشکل به‌ویژه در شرایطی که داده‌های آزمون دارای اغتشاش یا تنوع بالا هستند، شدت بیشتری پیدا می‌کند.  

برای مقابله با این مسئله، دسته‌ای از روش‌ها به طور خاص برای کنترل روند انطباق معرفی شده‌اند که هدف اصلی آن‌ها ایجاد تعادل میان سازگاری با داده‌های جدید و حفظ دانش منبع است. این روش‌ها معمولاً از سازوکار‌هایی همچون بازگردانی پارامترها، محدودسازی به‌روزرسانی‌ها بر اساس اهمیت پارامترها و انتقال تدریجی دانش بهره می‌گیرند.  

به طور خاص، معماری‌های معلم–دانش‌آموز\پانویس{Teacher–Student} نقش مهمی در انتقال تدریجی دانش ایفا می‌کنند. در این چارچوب، مدل معلم مسئول تثبیت و نگهداری دانش پیشین است، در حالی که مدل دانش‌آموز وظیفه‌ی انطباق با داده‌های جدید را بر عهده دارد. هرچند این رویکرد می‌تواند به شکل مؤثری از فراموشی جلوگیری کند، اما معمولاً موجب کندی فرایند انطباق می‌شود، زیرا نیازمند چندین گذر رو به جلو و تبادل مداوم اطلاعات میان دو مدل است.  

در مجموع، روش‌های ضد فراموشی می‌کوشند تا بدون افت محسوس در کیفیت بازنمایی‌های پیشین، امکان سازگاری مدل با توزیع‌های جدید داده را فراهم کنند. تمرکز اصلی این دسته از رویکردها بر کنترل روند انطباق و نه لزوماً بر ریشه‌یابی یا حذف عوامل ایجاد خطا می‌باشد.

\زیرزیرقسمت{روش \lr{COTTA}} 

روش \lr{CoTTA}\پانویس{Continual Test-Time Adaptation} با هدف سازگاری تدریجی و پیوسته در شرایط آزمون طراحی شده است. ایده‌ی اصلی این روش بر پایه‌ی استفاده از سازوکارهای به‌روزرسانی نرم و جلوگیری از تغییرات ناگهانی در پارامترهای شبکه است تا هم مشکل انباشت خطا کاهش یابد و هم از فراموشی فاجعه‌بار جلوگیری شود.  

\lr{CoTTA} از سه مؤلفه‌ی کلیدی بهره می‌برد:  
\begin{itemize}
	\item \textbf{میانگین‌گیری نمایی مدل}\پانویس{Exponential Moving Average (EMA)}: یک نسخه‌ی میانگین‌گیری‌شده از پارامترهای مدل به عنوان مدل معلم نگهداری می‌شود تا ثبات و دانش قبلی را تضمین کند.  
	\item \textbf{تشدید داده‌ها}\پانویس{Stochastic Augmentation}: هر نمونه‌ی ورودی با چندین تبدیل تصادفی پردازش می‌شود تا تنوع بیشتری ایجاد شود و مدل از بیش‌اعتماد شدن روی یک نمایش خاص جلوگیری کند. این ایده الهام گرفته شده از مشاهداتی است که نشان می‌دهد مدل‌های میانگین‌گیری وزنی در طول مراحل آموزش اغلب عملکرد دقیق‌تری نسبت به مدل نهایی ارائه می‌دهند \cite{polyak1992, tarvainen2017}.
	
	\item \textbf{بازگردانی وزنی}\پانویس{Weight Restoration}: برای جلوگیری از انحراف شدید، بخشی از پارامترها به صورت تصادفی به مقدار اولیه‌ی آموزش منبع بازگردانده می‌شوند. این سازوکار نوعی حافظه‌ی پارامترها ایجاد می‌کند که جلوی فراموشی سریع دانش منبع را می‌گیرد.  
\end{itemize}

تابع هزینه‌ی اصلی \lr{CoTTA} مبتنی بر یک سازوکار خودآموزی\پانویس{Self-Training} است که پیش‌بینی‌های مدل معلم به عنوان برچسب‌ برای مدل دانش‌آموز طبق رابطه‌ی \ref{eq:cotta} در نظر گرفته می‌شوند\cite{cotta}:  

\begin{equation}
	L_{\text{CoTTA}} = \frac{1}{K} \sum_{k=1}^{K} H(f_{\theta}(a_k(x)), f_{\theta'}(x)),
	\label{eq:cotta}
\end{equation}

که در آن \(a_k(x)\) نمایش‌های مختلف نمونه‌ی ورودی \(x\) تحت تبدیل‌های تصادفی هستند، \(f_{\theta}\) مدل دانش‌آموز و \(f_{\theta'}\) مدل معلم را نشان می‌دهد و \(H(\cdot, \cdot)\) آنتروپی متقابل بین دو توزیع است.  

سازوکار بازگردانی وزنی به صورت زیر طبق رابطه‌ی \ref{eq:update} تعریف می‌شود\cite{cotta}:  

\begin{equation}
	\theta \leftarrow m \cdot \theta + (1-m) \cdot \theta_0
	\label{eq:update}
\end{equation}


که در آن \(\theta_0\) پارامترهای اولیه‌ی مدل منبع، و \(m\) یک ضریب بازگردانی کوچک است. این به‌روزرسانی مانع انحراف بیش از حد مدل دانش‌آموز از دانش منبع می‌شود.  

در مجموع، \lr{CoTTA} با ترکیب سه سازوکار میانگین‌گیری نمایی مدل، تشدید داده و بازگردانی وزنی، توانسته است تعادلی میان سازگاری با داده‌های جدید و جلوگیری از فراموشی فاجعه‌بار برقرار کند. هرچند این روش نسبت به رویکردهای ساده‌تر هزینه‌ی محاسباتی بیشتری دارد، اما در محیط‌های پویا و ناشناخته عملکرد پایدارتر و قابل اعتمادتری ارائه می‌دهد.  

\زیرزیرقسمت{روش \lr{EATA}} 

روش \lr{EATA}\پانویس{Efficient Anti-forgetting Test-time Adaptation} با هدف افزایش کارایی در انطباق زمان آزمون و جلوگیری از پدیده‌ی فراموشی فاجعه‌بار توسعه یافته است. این روش بر پایه‌ی دو مؤلفه‌ی اصلی انتخاب کارآمد نمونه‌ها برای کاهش محاسبات غیرضروری و منظم‌ساز ضدفراموشی برای حفظ دانش درون‌توزیع می‌باشد. 

در بخش انتخاب نمونه‌ها، ایده‌ی اصلی آن است که همه‌ی داده‌های ورودی برای انطباق به‌کار گرفته نشوند؛ بلکه تنها نمونه‌های معتبر و غیرتکراری در فرآیند {پس‌انتشار\پانویس{Backpropagation}} نقش داشته باشند. برای این منظور، به هر نمونه یک وزن تطبیقی \(S(x)\) نسبت داده می‌شود که ترکیبی از دو معیار است:  

\textbf{۱. قابلیت اعتماد (اعتمادپذیری نمونه):} این معیار بر اساس آنتروپی پیش‌بینی تعریف می‌شود. طبق رابطه‌ی \ref{eq:eata-sent}، اگر آنتروپی نمونه کمتر از یک آستانه‌ی از پیش تعیین‌شده \(E_0\) باشد، آن نمونه به‌عنوان قابل اعتماد در نظر گرفته می‌شود\cite{eata}:  

\begin{equation}
	S_{\text{ent}}(x) = \frac{1}{\exp(E(x; \Theta) - E_0)} \cdot \mathbb{I}_{\{E(x; \Theta) < E_0\}}.
	\label{eq:eata-sent}
\end{equation}

\textbf{۲. نبود افزونگی (تنوع نمونه‌ها):} برای جلوگیری استفاده‌ی مکرر از نمونه‌های مشابه، از یک میانگین متحرک نمایی برای ردیابی خروجی‌های مدل روی داده‌های آزمون گذشته استفاده می‌شود. این میانگین طبق رابطه‌ی \ref{eq:eata-moving-average} به‌روزرسانی می‌شود\cite{eata}:  

\begin{equation}
	m_t = \alpha \bar{y}_t + (1 - \alpha) m_{t-1}, \quad t \geq 1,
	\label{eq:eata-moving-average}
\end{equation}

که در آن \(\bar{y}_t = \frac{1}{n} \sum_{k=1}^n \hat{y}^k_t\) میانگین پیش‌بینی یک دسته شامل \(n\) نمونه در تکرار \(t\) است. سپس برای هر نمونه‌ی جدید \(x\)، شباهت کسینوسی بین \(f_{\Theta}(x)\) و \(m_{t-1}\) محاسبه شده و طبق رابطه‌ی \ref{eq:eata-sdiv} اگر این شباهت کمتر از آستانه‌ی \(\epsilon\) باشد، نمونه غیرتکراری تلقی می‌شود\cite{eata}:  

\begin{equation}
	S_{\text{div}}(x) = \mathbb{I}_{\{\cos(f_{\Theta}(x), m_{t-1}) < \epsilon\}}.
	\label{eq:eata-sdiv}
\end{equation}

وزن کلی نمونه بر اساس ترکیب این دو معیار و مطابق رابطه‌ی \ref{eq:eata-s-total} تعریف می‌شود\cite{eata}:  

\begin{equation}
	S(x) = S_{\text{ent}}(x) \cdot S_{\text{div}}(x).
	\label{eq:eata-s-total}
\end{equation}

تنها نمونه‌هایی با \(S(x) > 0\) در فرآیند به‌روزرسانی دخالت داده می‌شوند. در نتیجه، تابع هزینه‌ی انطباق طبق رابطه‌ی \ref{eq:eata-entropy} به شکل زیر نوشته می‌شود\cite{eata}:  

\begin{equation}
	\min_{\tilde{\Theta}} S(x) E(x; \Theta) = - S(x) \sum_{y \in C} f_{\Theta}(y|x) \log f_{\Theta}(y|x),
	\label{eq:eata-entropy}
\end{equation}

که در آن \(C\) فضای خروجی مدل است و تنها پارامترهای مربوط به لایه‌های نرمال‌سازی دسته‌ای \(\tilde{\Theta}\) به‌روزرسانی می‌شوند.  

برای مقابله با فراموشی فاجعه‌بار، \lr{EATA} یک منظم‌ساز وزنی مبتنی بر اطلاعات فیشر معرفی می‌کند که مانع تغییر بیش‌ازحد پارامترهای مهم مدل می‌شود. این منظم‌ساز طبق رابطه‌ی \ref{eq:eata-regularizer} تعریف می‌گردد\cite{eata}:  

\begin{equation}
	R(\tilde{\Theta}; \tilde{\Theta}_0) = \sum_{\theta_i \in \tilde{\Theta}} \omega(\theta_i) (\theta_i - \theta_i^0)^2,
	\label{eq:eata-regularizer}
\end{equation}

که در آن \(\tilde{\Theta}_0\) وزن‌های اولیه مدل و \(\omega(\theta_i)\) اهمیت پارامتر \(\theta_i\) است. اهمیت پارامترها بر اساس ماتریس فیشر تخمین زده می‌شود. از آنجا که داده‌های آموزشی برچسب‌دار در دسترس نیستند، مجموعه‌ای از نمونه‌های درون‌توزیع بدون برچسب جمع‌آوری شده و با استفاده از مدل اولیه شبه‌برچسب‌های سخت برای آن‌ها تولید می‌شود. سپس با استفاده از این داده‌ی شبه‌برچسب‌دار، اهمیت وزنی محاسبه می‌گردد.  

در نهایت، تابع هدف کلی روش \lr{EATA} طبق رابطه‌ی \ref{eq:eata-total} به صورت زیر خواهد بود\cite{eata}:  

\begin{equation}
	\min_{\tilde{\Theta}} S(x) E(x; \Theta) + \beta R(\tilde{\Theta}; \tilde{\Theta}_0),
	\label{eq:eata-total}
\end{equation}

که در آن \(\beta\) پارامتر موازنه بین انطباق و جلوگیری از فراموشی است.  

به‌طور خلاصه، \lr{EATA} یک روش برای حذف نمونه‌های تکراری ارائه می‌دهد که سرعت انطباق را افزایش داده و همین‌طور خود این امر می‌تواند به جلوگیری از انباشت خطا و فراموشی فاجعه‌بار کمک کند. این روش تنها نیازمند چند گام پس‌انتشار روی نمونه‌های منتخب است و سربار محاسباتی زیادی ندارد. با این حال، همچنان برای انتخاب نمونه‌های مطمئن بر آنتروپی تکیه دارد که همان‌طور که پیش‌تر اشاره شد و در فصل آینده نیز بررسی خواهد شد، لزوماً معیار قابل اعتمادی برای حذف نمونه‌های دارای اغتشاش نبوده و حتی می‌تواند در مواردی موجب تسریع فروپاشی مدل شود.  

\newpage

\قسمت{جمع‌بندی}

در این فصل ابتدا به بررسی مبدل‌های بینایی و لایه‌های نرمال‌ساز در شبکه‌های عصبی پرداختیم و اهمیت آن‌ها را در مسئله‌ی انطباق زمان آزمون بیان کردیم. همچنین تفاوت میان انواع لایه‌های نرمال‌سازی و نقش آن‌ها در پایداری و کیفیت انطباق بررسی شد. در ادامه، طیفی از روش‌های مختلف انطباق زمان آزمون برخط معرفی گردید و مزایا و چالش‌های هر یک مورد بحث قرار گرفت. جمع‌بندی کلی این روش‌ها نیز در جدول \ref{table:tta-comparison} ارائه شده‌است.

با مرور پژوهش‌های پیشین مشاهده می‌شود که همچنان چالش‌های متعددی در این حوزه وجود دارد. بسیاری از روش‌ها به مجموعه‌ای از فراپارامترهای متعدد وابسته‌اند که پیاده‌سازی آن‌ها را در کاربردهای عملی با چالش‌های جدی روبه‌رو می‌کند، به‌ویژه آنکه برخی روش‌ها نسبت به مقدار این فراپارامترها حساسیت بالایی دارند. علاوه بر این، سازوکار‌های موجود برای حذف نمونه‌های نامطمئن عمدتاً بر مبنای معیارهایی همچون آنتروپی یا بیشینه‌ی خروجی مدل تعریف شده‌اند؛ حال آنکه این معیارها در مواجهه با داده‌های دارای اغتشاش یا خارج از توزیع می‌توانند کاملاً گمراه‌کننده باشند.  

در فصل بعد، تلاش خواهیم کرد رویکردی متناسب با الزامات انطباق زمان آزمون برخط ارائه دهیم؛ رویکردی که از یک‌سو سادگی پیاده‌سازی و حساسیت کم به فراپارامترها را داشته باشد و از سوی دیگر، بخشی از مشکلات مطرح‌شده در پژوهش‌های پیشین را برطرف کند.

\begin{table}[h!]
	\centering
	\small
	\caption{جمع‌بندی روش‌های مختلف انطباق در زمان آزمون برخط}
	\label{table:tta-comparison}
	\begin{tabular}{p{2.5cm} p{6cm} p{6cm}}
		\toprule
		\textbf{روش} & \textbf{مزایا} & \textbf{معایب} \\ 
		\midrule
		\lr{TENT}\cite{ftent} & پیاده‌سازی ساده؛ بدون نیاز به پارامترهای اضافی. & مستعد راه‌حل‌های بدیهی؛ استفاده مستقیم از آنتروپی و بهینه‌ساز آزاد بدون محدودیت در به‌روزرسانی وزن‌های مدل، می‌تواند منجر به فروپاشی مدل به‌ویژه در حالتهای پیوسته شود. \\ 
		\midrule
		\lr{SAR}\cite{sar} & تحلیل دقیق موارد شکست با توجه به تغییرات ناگهانی در اندازه‌ی گرادیان‌ها و طراحی تابع هدف بر اساس این مشاهدات که به طور موثر از فروپاشی مدل در سنارو‌های مختلف جلوگیری می‌کند.& نیازمند تنظیم دقیق لایه‌هایی که باید منجمد یا قابل آموزش باشند، زیرا لایه‌های پایانی معمولاً به تغییرات حساس‌اند؛ تعداد فرا‌پارامتر‌های روش زیاد است و لزوم انتخاب عمق، کران و شعاع اغتشاش مناسب، استفاده از این روش را در محیط‌های عملیاتی و برخط دچار چالش جدی می‌کند. \\ 
		\midrule
		\lr{DeYo}\cite{deyo} & تنها نمونه‌هایی که آنتروپی پایین و اختلاف معنی‌دار دارند وارد انطباق می‌شوند؛ تابع هزینه تمرکز روی نمونه‌های معنادار را تضمین می‌کند و از فروپاشی مدل جلوگیری می‌کند. & سربار محاسباتی ناشی از اعمال آشفتگی؛ نیاز به تنظیم دقیق آستانه‌ها؛ اگر مدل بر اساس ویژگی‌های محلی یک نمونه پیش‌بینی دقیقی ارائه دهد، اعمال آشفتگی پیش‌بینی مدل را تغییر نمی‌دهد و در این موارد این نمونه‌های مطمئن از فرایند انطباق حذف می‌شوند. \\ 
		\bottomrule
	\end{tabular}
\end{table}

\begin{table}[h!]\ContinuedFloat
	\centering
	\small
	\caption{جمع‌بندی روش‌های مختلف انطباق در زمان آزمون برخط}
	\begin{tabular}{p{2.5cm} p{6cm} p{6cm}}
		\toprule
		\textbf{روش} & \textbf{مزایا} & \textbf{معایب} \\ 
		\midrule
		\lr{TEA}\cite{tea} & استفاده از انرژی به جای آنتروپی؛ ترجیح عدم قطعیت به جای بیش‌اعتمادی؛ موثر زمانی که مدل ذاتاً نامطمئن است؛ کاهش خطای \lr{ECE} و \lr{MCE}. & کنترل مستقیمی روی تابع انرژی وجود ندارد؛  مشابه آنتروپی استفاده از تابع انرژي بدون اعمال ساز‌وکار‌های مناسب تشخیص و حذف نمونه‌های نامطمئن مدل را مستعد فروپاشی می‌کند؛ استفاده از \lr{SGLD} باعث کندی زیاد شده و استفاده از این روش را در محیط‌های بلادرنگ عملا ناممکن می‌کند. \\ 
		\midrule
		\lr{AEA}\cite{aea} & ترکیب انرژی و آنتروپی باعث تسریع انطباق و تعادل میان بیش‌اعتمادی و عدم قطعیت می‌شود؛ عملکرد بهتر در برخی مجموعه‌داده‌ها و مدل‌ها مشاهده شده است. & حساسیت بالا نسبت به هم‌ترازی  توابع انرژی و آنتروپی در تابع هزینه؛ نیاز به تنظیم دقیق پارامترهای $\lambda_1$ و $\lambda_2$ که در شرایط برخط بسیار چالش‌برانگیز است. \\ 
		\midrule
		\lr{T3A}\cite{t3a} & استفاده از مجموعه پشتیبانی و فیلتر آنتروپی باعث پایدار ماندن پیش‌بینی‌ها و جلوگیری از تاثیر نمونه‌های با برچسب موقت نادرست می‌شود. & وابستگی به انتخاب درست آستانه آنتروپی \(\alpha_k\)؛ احتمال اختصاص اشتباه برچسب موقت که عملکرد مدل را کاهش می‌دهد. \\ 
		\midrule
		\lr{SHOT}\cite{shot} & ترکیب مزایای روش‌های مبتنی بر آنتروپی با سازوکار پایدار برچسب‌زنی خودکار؛ انطباق سریع و کاهش اثر برچسب‌های دارای اغتشاش. & برای مسائل برخط و جریان داده‌ای طراحی نشده است و به کار گیری آن برای این دسته مسائل نیاز به ترکیب کردن آن با روش‌های دیگر دارد؛ راهکار مشخصی برای جلوگیری از فروپاشی مدل ارائه نمی‌کند؛ حساس به عدم توازن برچسب‌ها در نمونه‌ها. \\ 
		\bottomrule
	\end{tabular}
\end{table}

\begin{table}[h!]\ContinuedFloat
	\centering
	\small
	\caption{جمع‌بندی روش‌های مختلف انطباق در زمان آزمون برخط}
	\begin{tabular}{p{2.5cm} p{6cm} p{6cm}}
		\toprule
		\textbf{روش} & \textbf{مزایا} & \textbf{معایب} \\ 
		\midrule
		\lr{REM}\cite{rem} & حفظ استحکام و انعطاف‌پذیری مدل در زمان آزمون؛ کاهش حساسیت به تغییرات ناگهانی داده‌ها. & طول زنجیره و نسبت پوشش‌ها به عنوان فراپارامترهای کلیدی نیازمند تنظیم دقیق هستند؛ سازوکار مشخصی برای حذف نمونه‌های نامطمئن ارائه نمی‌دهد که در حالتهای برخط چالش‌برانگیز است. \\ 
		\midrule
		\lr{PETAL}\cite{petal} & کاهش اثر فراموشی فاجعه‌بار و انباشت خطا؛ انطباق پایدار روی داده‌های بدون برچسب؛ عدم نیاز به دسترسی به داده‌های منبع و حفظ دانش اولیه مدل. & پیچیدگی محاسباتی بالا؛ حساسیت به تنظیم فراپارامترها؛ محدودیت در جریان داده‌های بسیار تغییرپذیر؛ تقریب گوسی ممکن است تمامی پیچیدگی فضای وزن‌ها را مدل‌سازی نکند. \\ 
		\midrule
		\lr{COTTA}\cite{cotta} & جلوگیری از فروپاشی مدل با استفاده از معماری معلم-دانش‌آموز و میانگین‌گیری متحرک؛ ساده و مؤثر در کاهش انباشت خطا. & از نظر محاسباتی بسیار پرهزینه؛ انطباق کند؛ کاهش انعطاف ‌پذیری. \\ 
		\midrule
		\lr{EATA}\cite{eata} & ارائه روشی برای حذف نمونه‌های تکراری که سرعت انطباق را افزایش می‌دهد و به کاهش انباشت خطا و فراموشی فاجعه‌بار کمک می‌کند؛ تنها نیازمند چند گام پس‌انتشار روی نمونه‌های منتخب بدون سربار محاسباتی زیاد. & تکیه بر آنتروپی برای انتخاب نمونه‌های مطمئن که لزوماً معیار قابل اعتمادی برای حذف نمونه‌های دارای اغتشاش نیست و می‌تواند در مواردی فروپاشی مدل را تسریع کند. \\ 
		\bottomrule
	\end{tabular}
\end{table}